\projectprogress
\label{sec:progress}

\subsection{Pipeline Architecture}
The project is implemented as several Python modules that will all be combined
in a pipeline. Four sequential stages are in scope for this interim report; the
fifth stage (neural-network parameter optimization) is deferred to the second
half of the semester. Each stage is also exposed as a standalone module so that
any subset of the workflow can be re-executed in isolation. The current state of
completion of each stage is summarized in Table~\ref{tab:pipeline_status}.

\begin{table}[h]
	\centering
	\caption{Pipeline stages and current implementation status.}
	\label{tab:pipeline_status}
	\renewcommand{\arraystretch}{1.2}
	\begin{tabular}{|c|p{6.5cm}|c|c|}
		\hline
		\textbf{Stage} & \textbf{Description} & \textbf{Framework} & \textbf{Status} \\ \hline
		1 & Random alloy configuration generation     & Python                       & Complete \\ \hline
		2 & MD elastic-property evaluation             & LAMMPS Python API            & Complete \\ \hline
		3 & DFT reference calculations                  & Quantum Espresso + ASE       & Complete \\ \hline
		4 & MEAM initialization and library merging     & Python                       & Complete \\ \hline
	\end{tabular}
\end{table}

The data flow between stages is shown schematically in
Figure~\ref{fig:pipeline}. Stage~1 emits a JSON description of an alloy
configuration; Stage~2 simulates that configuration in LAMMPS and writes the
measured elastic moduli to a shared results database; Stage~3 produces DFT
reference values for every selected element and binary pair and stores them in
a reference-data database; Stage~4 reads both of these and emits a merged,
DFT-initialized pair of MEAM files. A divergent composition-only neural-network
branch (dashed in Figure~\ref{fig:pipeline}) consumes the Stage~2 elastic
database as a sanity check but never feeds back into the potential-fitting
flow; it is described in Section~\ref{sec:comp_nn}.

\begin{figure}[h]
	\centering
	\begin{tikzpicture}[
		stage/.style={rectangle, rounded corners, draw=black, fill=gray!10,
		              text width=2.6cm, minimum height=1.0cm, align=center, font=\small\bfseries},
		aux/.style={rectangle, rounded corners, draw=black, dashed, fill=blue!5,
		              text width=2.6cm, minimum height=1.0cm, align=center, font=\small\bfseries},
		arrow/.style={-Stealth, thick},
		auxarrow/.style={-Stealth, thick, dashed},
		label/.style={font=\scriptsize, align=center},
		node distance=0.4cm and 0.4cm,
	]
		\node[stage] (s1) {Stage 1\\Configurations};
		\node[stage, right=of s1] (s2) {Stage 2\\MD: $E,\nu$};
		\node[stage, right=of s2] (s3) {Stage 3\\DFT};
		\node[stage, right=of s3] (s4) {Stage 4\\MEAM init};

		\draw[arrow] (s1) -- (s2);
		\draw[arrow] (s2) -- (s3);
		\draw[arrow] (s3) -- (s4);

		\node[label, below=0.1cm of s1] {composition\\records};
		\node[label, below=0.1cm of s2] {elastic\\database};
		\node[label, below=0.1cm of s3] {DFT\\reference};
		\node[label, below=0.1cm of s4] {merged MEAM\\library};

		% Divergent ML sanity-check branch: consumes the elastic database
		% from Stage 2 but does not feed back into the main pipeline.
		\node[aux, above=0.9cm of s3] (ml) {ML branch\\(\S\ref{sec:comp_nn})};
		\draw[auxarrow] (s2.north) |- (ml.west);
		\node[label, right=0.1cm of ml] {sanity check\\$(E,\nu)$ predictor};
	\end{tikzpicture}
	\caption{Four-stage pipeline architecture (solid arrows) plus the
	divergent ML sanity-check branch (dashed) described in
	Section~\ref{sec:comp_nn}. Each main stage is implemented as
	an independent Python module and communicates with the next one
	through JSON or MEAM-format files written to disk; the ML branch
	consumes the Stage~2 elastic database but never feeds back into
	the potential-fitting flow.}
	\label{fig:pipeline}
\end{figure}

\newpage
